\documentclass[../main.tex]{subfiles}
\graphicspath{{\subfix{../image/}}}

\begin{document}

\chapter{向量丛}\label{chapter:ch10-向量丛}

在第~\cref{chapter:ch03}~中，我们看到光滑流形的切丛本身具有作为光滑流形的自然结构。我们在 $TM$ 上构造的自然坐标使其局部看起来像是 $M$ 的一个开子集与 $\mathbb{R}^n$ 的笛卡尔积。这种结构相当频繁地出现——为 $M$ 中每一点配备一个向量空间，并以局部看似 $M$ 与 $\mathbb{R}^n$ 的笛卡尔积、但整体可能被“扭曲”的方式粘合在一起。这样的结构称为\textbf{向量丛}，是本章的主要课题。

本章从向量丛的定义开始，并描述若干例子。最值得注意的例子当然是光滑流形的切丛。接着我们讨论向量丛的局部截面与整体截面（在切丛的情形中对应于向量场）。本章继续讨论丛之间映射的自然概念，即\textbf{丛同态}，以及向量丛自身也是向量丛的子集，即\textbf{子丛}。本章末尾，我们简要介绍向量丛的一个重要推广，称为\textbf{纤维丛}。

关于流形上向量丛与纤维丛，存在着深刻而广博的理论，我们甚至无法触及。我们引入它们主要是为了获得一套方便的语言来谈论切丛及其类似结构；正如你将在接下来几章中看到的，这类结构在光滑流形上大量存在。

\subfile{Chapter-10/section-01.tex}
\subfile{Chapter-10/section-02.tex}
\subfile{Chapter-10/section-03.tex}
\subfile{Chapter-10/section-04.tex}
\subfile{Chapter-10/section-05.tex}
\end{document}